%% 2-15 比例と反比例（第11回で書き足し）
\makeatletter
% ばね（たて）  \bkiipspring{(x,上のはし)}{全体の長さ}  コイルの上下に短いまっすぐな部分をつける
\newcommand\bkiipspring[2]{\begin{scope}[shift={#1}]
  \draw[bkGray!90!black, line width=0.7pt] (0,0) -- (0,-2.5);
  \draw[bkGray!90!black, line width=0.7pt, decorate, decoration={coil, aspect=0.45, segment length=1.4mm, amplitude=2.2mm}] (0,-2.5) -- (0,{-(#2)+2.5});
  \draw[bkGray!90!black, line width=0.7pt] (0,{-(#2)+2.5}) -- (0,{-(#2)});
\end{scope}}
% おもり  \bkiipweight{(x,上のはし)}  はば7mm・高さ6mm
\newcommand\bkiipweight[1]{\begin{scope}[shift={#1}]
  \filldraw[fill=bkGray!45, draw=bkGray!90!black, line width=0.6pt, rounded corners=0.8mm] (-3.5,-6) rectangle (3.5,0);
  \fill[white, opacity=0.35] (-2.6,-1.6) rectangle (-1.4,-4.8);
\end{scope}}
% 式を、まん中に1行で置く  \bkiipline{…}
\newcommand\bkiipline[1]{\par\addvspace{3pt}{\centering #1\par}\addvspace{3pt}}
% 天井の板  \bkiipceil{(x1,y)}{(x2,y)}
\newcommand\bkiipceil[2]{\fill[bkGray!35] #1 rectangle ($#2+(0,2.2)$);
  \draw[bkGray!90!black, line width=0.8pt] #1 -- #2;}
% 2行の表に「何倍」の矢印をつけた図
%   \bkiiptable{上の行の名前}{上の数のリスト}{下の行の名前}{下の数のリスト}
%   矢印は \bkiiparrow{上か下(1か-1)}{左の列}{右の列}{文字}{色}  列は0から
\newcommand\bkiiptable[4]{%
  \foreach \i in {0,...,6} \draw[bkLine, line width=0.6pt] (32+\i*13,-6) -- (32+\i*13,6);
  \draw[bkLine, line width=0.6pt] (0,6) rectangle (110,-6) (0,0) -- (110,0);
  \fill[bkTheme!12] (0.3,5.7) rectangle (31.7,-5.7);
  \node[font=\sffamily\small, anchor=west] at (2,3) {#1};
  \node[font=\sffamily\small, anchor=west] at (2,-3) {#3};
  \foreach \v [count=\i from 0] in {#2} \node[font=\small] at (38.5+\i*13,3) {$\v$};
  \foreach \v [count=\i from 0] in {#4} \node[font=\small] at (38.5+\i*13,-3) {$\v$};}
\newcommand\bkiiparrow[5]{%
  \pgfmathsetmacro\bkiip@a{38.5+#2*13}\pgfmathsetmacro\bkiip@b{38.5+#3*13}%
  \draw[-{Stealth[length=4pt]}, #5, line width=0.9pt, shorten <=1mm, shorten >=1mm] (\bkiip@a,#1*6.8) .. controls (\bkiip@a+3,#1*11.5) and (\bkiip@b-3,#1*11.5) .. (\bkiip@b,#1*6.8);
  \ifnum#1>0 \def\bkiip@anc{south}\else\def\bkiip@anc{north}\fi
  \node[font=\sffamily\footnotesize, text=#5!80!black, anchor=\bkiip@anc, inner sep=0.5pt] at ({(\bkiip@a+\bkiip@b)/2},#1*10.8) {#4};}
\makeatother

\bkchapter{2}{15}{\R{比例}{ひ|れい}と\R{反比例}{はん|ぴ|れい}}{\bklevelrange{1}{2}}{\R{比例}{ひ|れい}と\R{反比例}{はん|ぴ|れい}}
\bkmichishirube{2-3,2-8,2-10,2-14}{3-3,3-4,3-7,5-12,6-6}

\begin{goal}
\begin{itemize}[label={\color{bkTheme}●}, leftmargin=1.6\zw, itemsep=0pt, topsep=1pt]
\item \bklevel{1}\enspace \R{表}{ひょう}から「$x$ が2\R{倍}{ばい}になると $y$ も2\R{倍}{ばい}」を\R{見}{み}つけて、$y=ax$ の\R{式}{しき}とグラフに\R{表}{あらわ}せる。
\item \bklevel{2}\enspace ばねの\R{力}{ちから} $F=kx$ と\R{電気}{でん|き}の\R{式}{しき} $V=RI$ を、\R{表}{ひょう}から\R{作}{つく}って\R{使}{つか}える。
\item \bklevel{2}\enspace \R{反比例}{はん|ぴ|れい}（$xy$ が\R{一定}{いっ|てい}）を\R{見分}{み|わ}け、「2\R{倍}{ばい}にするとどうなるか」を\R{答}{こた}えられる。
\end{itemize}
\end{goal}

\begin{hatena}
\noindent\begin{minipage}[c]{0.62\linewidth}
ばねにおもりを1\R{個}{こ}つるすと、\qty{2}{cm}のびました。おもりを2\R{個}{こ}にすると、\R{何}{なん}cmのびるかな？　10\R{個}{こ}つるしたら？

ばねを2\R{倍}{ばい}のばすには、2\R{倍}{ばい}の\R{力}{ちから}がいるのでしょうか。それとも、もっと\R{大}{おお}きな\R{力}{ちから}がいるのでしょうか。
\end{minipage}\hfill
\begin{minipage}[c]{0.34\linewidth}\centering
\begin{tikzpicture}[x=1mm, y=1mm]
  \bkiipceil{(-1,0)}{(42,0)}
  \draw[bkRed!70, line width=0.5pt, densely dotted] (1,-20) -- (26,-20);
  \bkiipspring{(4,0)}{20}
  \bkiipspring{(18,0)}{28}\bkiipweight{(18,-28)}
  \draw[bkGray!70, line width=0.7pt, dashed] (37,0) -- (37,-36);
  \node[font=\sffamily\bfseries\Large, text=bkOrange!80!black] at (37,-25) {?};
  \node[bk tick, anchor=north] at (4,-38) {0\R{個}{こ}};
  \node[bk tick, anchor=north] at (18,-38) {1\R{個}{こ}};
  \node[bk tick, anchor=north] at (37,-38) {2\R{個}{こ}};
  \draw[{Stealth[length=3pt]}-{Stealth[length=3pt]}, bkRed, line width=0.6pt] (24,-20) -- (24,-28);
  \node[font=\sffamily\scriptsize, text=bkRed!80!black, anchor=west] at (24.5,-31.5) {2\,cm};
\end{tikzpicture}
\end{minipage}
\end{hatena}

\clearpage
%% ---------------- やってみよう・思い出そう ----------------
\begin{yattemiyou}[ばねののびをはかろう]
\noindent\begin{minipage}[c]{0.6\linewidth}
\begin{enumerate}[label={\bkcircled[bkGreen]{\arabic*}}, leftmargin=2.4\zw, itemsep=1pt, topsep=0pt]
\item ばねをスタンドにつるし、\R{横}{よこ}に\R{立}{た}てたものさしで、ばねの\R{下}{した}のはしの\R{目}{め}もりを\R{読}{よ}もう。
\item 1\R{個}{こ}\qty{50}{g}のおもりを1\R{個}{こ}ずつふやしてつるし、そのたびに\R{目}{め}もりを\R{読}{よ}もう。
\item \R{読}{よ}んだ\R{目}{め}もりから、はじめの\R{目}{め}もりをひいて、のびを\R{表}{ひょう}にしよう（あと－まえ、2-3）。
\end{enumerate}
\end{minipage}\hfill
\begin{minipage}[c]{0.36\linewidth}\centering
\begin{tikzpicture}[x=1mm, y=1mm]
  % スタンド
  \fill[bkGray!60] (-6,-50) rectangle (26,-47.5);
  \fill[bkGray!70] (-4,-47.5) rectangle (-2,6);
  \fill[bkGray!70] (-4,4) rectangle (14,6);
  \bkiipspring{(10,4)}{24}
  % ばねの下のはしにつけた目じるし（ここで目もりを読む）
  \fill[bkRed] (10,-20) -- (14,-19) -- (14,-21) -- cycle;
  \draw[bkGray!90!black, line width=0.7pt] (10,-20) -- (10,-21.5);
  \bkiipweight{(10,-21.5)}\bkiipweight{(10,-28)}
  % ものさし（台の上に立てる）
  \filldraw[fill=bkYellowBg, draw=bkYellow!80!black, line width=0.6pt] (19,-47.5) rectangle (24,2);
  \foreach \k in {0,...,9} \draw[bkInk, line width=0.4pt] (19,{-45+\k*5}) -- (21,{-45+\k*5});
  \draw[bkRed, line width=0.6pt, densely dashed] (14.5,-20) -- (19,-20);
\end{tikzpicture}
\end{minipage}
\par\vspace{4pt}
ミナさんがはかった\R{結果}{けっ|か}です。
\begin{bkcenter}
\begin{bktable}{colspec={Q[l,m,wd=30mm] *{6}{Q[c,m,wd=9.5mm]}}, rowsep=1.5pt}
おもりの\R{数}{かず}（\R{個}{こ}） & 0 & 1 & 2 & 3 & 4 & 5 \\
のび（cm） & 0 & 2.0 & 4.1 & 5.9 & 8.0 & 10.0 \\
\end{bktable}
\end{bkcenter}
のびは、おもりの\R{数}{かず}とどんな\R{関係}{かん|けい}になっているかな？
\end{yattemiyou}

\begin{omoidasou}
\bkrecall{2-3}{わり\R{算}{ざん}で「1あたりの\R{量}{りょう}」をつくる。\R{速}{はや}さ＝\R{道}{みち}のり÷\R{時間}{じ|かん}}
\bkrecall{2-6}{はかった\R{値}{あたい}にはずれがある。\R{何回}{なん|かい}かの\R{値}{あたい}をならして\R{考}{かんが}える。}
\bkrecall{2-8}{\R{振}{ふ}り\R{子}{こ}の\R{周期}{しゅう|き}は、\R{糸}{いと}の\R{長}{なが}さを4\R{倍}{ばい}にすると2\R{倍}{ばい}になる。}
\bkrecall{2-10}{\R{単位}{たん|い}をそろえる：$\qty{2}{cm}\div\qty{100}{cm/m}=\qty{0.02}{m}$（1 mあたり100 cm）}
\bkrecall{2-14}{$V=RI$ の\R{両辺}{りょう|へん}を $R$ でわると $I=\dfrac{V}{R}$}
\end{omoidasou}

\bkfukidashi{delta}{\R{表}{ひょう}の\R{数}{かず}を、\R{横}{よこ}に\R{見}{み}たり、たてに\R{見}{み}たりすると、きまりが\R{見}{み}つかるよ。}

\clearpage
%% ---------------- 数学の道具 ----------------
\bksection{doguu}{\R{数学}{すう|がく}の\R{道具}{どう|ぐ}}

\bksubsection{1}{2\R{倍}{ばい}、3\R{倍}{ばい}になると：\R{比例}{ひ|れい}}

\leavevmode\label{bk:2-15:hirei}やってみようの\R{表}{ひょう}を、たてに\R{見}{み}てみましょう。のびをおもりの\R{数}{かず}でわると、おもり1\R{個}{こ}あたりののびがわかります（2-3の「1あたりの\R{量}{りょう}」）。
\begin{bkcenter}
\begin{bktable}{colspec={Q[l,m,wd=40mm] *{5}{Q[c,m,wd=12mm]}}, rowsep=1pt, cells={font=\small}}
おもりの\R{数}{かず}（\R{個}{こ}） & 1 & 2 & 3 & 4 & 5 \\
のび（cm） & 2.0 & 4.1 & 5.9 & 8.0 & 10.0 \\
\SetCell{bg=bkYellow!30} 1\R{個}{こ}あたりののび（cm） & 2.0 & 2.05 & 1.97 & 2.0 & 2.0 \\
\end{bktable}
\bktabcaption{のび÷\R{個数}{こ|すう}は、どれも\R{約}{やく}2（1.97は $\dfrac{1}{100}$ の\R{位}{くらい}までのがい\R{数}{すう}）}
\end{bkcenter}
どれも\R{約}{やく}\qty{2}{cm}です。2.05や1.97のように\R{少}{すこ}しずれるのは、はかるときのずれです（2-6）。ならすと、おもり1\R{個}{こ}あたり\qty{2}{cm}のびると\R{考}{かんが}えられます。そこで、おもりの\R{数}{かず}を $x$ \R{個}{こ}、のびを $y$\,cm として、ずれをならした\R{表}{ひょう}を\R{作}{つく}ります。

\begin{bkfigure}
\begin{tikzpicture}[x=1mm, y=1mm]
  \bkiiptable{おもり $x$（\R{個}{こ}）}{0,1,2,3,4,5}{のび $y$（cm）}{0,2,4,6,8,10}
  \bkiiparrow{1}{1}{2}{2\R{倍}{ばい}}{bkBlue}
  \bkiiparrow{-1}{1}{2}{2\R{倍}{ばい}}{bkOrange}
  \foreach \sg/\anc/\col in {1/south/bkBlue, -1/north/bkOrange} {
    \draw[-{Stealth[length=4pt]}, \col, line width=0.9pt, shorten <=1mm, shorten >=1mm] (51.5,\sg*6.8) .. controls (54.5,\sg*25) and (74.5,\sg*25) .. (77.5,\sg*6.8);
    \node[font=\sffamily\footnotesize, text=\col!80!black, anchor=\anc, inner sep=0.5pt] at (64.5,\sg*21) {3\R{倍}{ばい}};}
\end{tikzpicture}
\bkfigcaption{$x$ が2\R{倍}{ばい}・3\R{倍}{ばい}になると、$y$ も2\R{倍}{ばい}・3\R{倍}{ばい}になる。}
\end{bkfigure}

$x$ が1から2へ2\R{倍}{ばい}になると、$y$ も2から4へ2\R{倍}{ばい}。$x$ が1から3へ3\R{倍}{ばい}になると、$y$ も2から6へ3\R{倍}{ばい}。このように、$x$ が2\R{倍}{ばい}、3\R{倍}{ばい}、…になると、$y$ も2\R{倍}{ばい}、3\R{倍}{ばい}、…になるとき、「$y$ は $x$ に\textbf{\R{比例}{ひ|れい}}する」といいます。\bkidx{ひれい}{\R{比例}{ひ|れい}}

\textbf{\R{式}{しき}にする}\quad 1\R{個}{こ}あたり\qty{2}{cm}のびるので、$x$ \R{個}{こ}なら、2 cmの $x$ こ\R{分}{ぶん}で $2\times x$（cm）のびます（1-12：1つ\R{分}{ぶん}×いくつ\R{分}{ぶん}）。だから $y=2x$ と\R{書}{か}けます。$x$ が0でないところでは、$y\div x$ はいつも2です。

\R{一般}{いっ|ぱん}に、\R{比例}{ひ|れい}は $y=ax$ と\R{書}{か}けます。$a$ は「$x$ が1のときの $y$」、つまり1あたりの\R{量}{りょう}で、\textbf{\R{比例定数}{ひ|れい|てい|すう}}といいます。$y=ax$ の\R{両辺}{りょう|へん}を $x$ でわると（2-14）、$a=\dfrac{y}{x}$ です。

\textbf{なぜ2\R{倍}{ばい}になると2\R{倍}{ばい}？}\quad $y=ax$ で $x$ を2\R{倍}{ばい}の $2x$ にすると、$y$ は $a\times2x=2\times ax$ で、もとの $ax$ の2\R{倍}{ばい}になります。$a$ がどんな\R{数}{かず}でも\R{同}{おな}じです。

\bkhatenaref の\R{答}{こた}え：2\R{個}{こ}なら $2\times2=4$ で\qty{4}{cm}、10\R{個}{こ}なら $2\times10=20$ で\qty{20}{cm}のびます。ただし、ばねがもとにもどらなくなるほど\R{強}{つよ}く\R{引}{ひ}くと、\R{比例}{ひ|れい}しなくなります。

\bksubsection{2}{\R{比例}{ひ|れい}のグラフは、0を\R{通}{とお}るまっすぐな\R{線}{せん}}

\leavevmode\label{bk:2-15:graph}2つの\R{量}{りょう}の\R{関係}{かん|けい}は、グラフにすると\R{一目}{ひと|め}でわかります。\R{横}{よこ}の\R{軸}{じく}に $x$、たての\R{軸}{じく}に $y$ をとった\R{方眼}{ほう|がん}に、\R{表}{ひょう}の $x$ と $y$ の\R{組}{くみ}を1つずつ\R{点}{てん}でかきこみます。

\noindent\begin{minipage}[t]{0.49\linewidth}\centering
\begin{tikzpicture}
\begin{axis}[bkplot, width=58mm, height=50mm, xmin=0, xmax=5.6, ymin=0, ymax=11.5, xtick={0,...,5}, ytick={0,2,...,10},
  xlabel={$x$（\R{個}{こ}）}, ylabel={$y$（cm）}]
  \addplot[bkBlue, dashed, line width=0.8pt] coordinates {(2,0) (2,4) (0,4)};
  \addplot[bkGreen!70!black, line width=1.1pt] coordinates {(3,6) (4,6) (4,8)};
  \addplot[only marks, mark=*, mark size=1.8pt, bkInk] coordinates {(0,0) (1,2) (3,6) (4,8) (5,10)};
  \addplot[only marks, mark=*, mark size=2.4pt, bkRed] coordinates {(2,4)};
  \node[font=\sffamily\footnotesize, anchor=south west, inner sep=1pt] at (axis cs:0.15,4.4) {$x=2$、$y=4$};
  \node[font=\sffamily\scriptsize, text=bkGreen!50!black, anchor=north] at (axis cs:3.5,5.7) {${}+1$};
  \node[font=\sffamily\scriptsize, text=bkGreen!50!black, anchor=west] at (axis cs:4.1,7) {${}+2$};
\end{axis}
\end{tikzpicture}
\bkfigcaption{\R{点}{てん}のかき\R{方}{かた}と、まっすぐになるわけ}
\end{minipage}\hfill
\begin{minipage}[t]{0.49\linewidth}\centering
\begin{tikzpicture}
\begin{axis}[bkplot, width=58mm, height=50mm, xmin=0, xmax=5.6, ymin=0, ymax=11.5, xtick={0,...,5}, ytick={0,2,...,10},
  xlabel={$x$（\R{個}{こ}）}, ylabel={$y$（cm）}]
  \addplot[bkcurve, domain=0:5.4] {2*x};
  \addplot[bkcurve2, domain=0:5.4] {x};
  \addplot[only marks, mark=*, mark size=1.7pt, bkInk] coordinates {(1,2.0) (2,4.1) (3,5.9) (4,8.0) (5,10.0)};
  \node[font=\sffamily\footnotesize, text=bkBlue!80!black, anchor=east] at (axis cs:3.9,9.6) {$y=2x$};
  \node[font=\sffamily\footnotesize, text=bkOrange!80!black, anchor=north west] at (axis cs:4.2,3.8) {$y=x$};
\end{axis}
\end{tikzpicture}
\bkfigcaption{はかった\R{点}{てん}と $y=2x$、$y=x$}
\end{minipage}

\vspace{4pt}
\textbf{\R{点}{てん}のかき\R{方}{かた}}\quad \R{左}{ひだり}のグラフの\R{赤}{あか}い\R{点}{てん}は、$x=2$、$y=4$ を\R{表}{あらわ}す\R{点}{てん}です。0から\R{右}{みぎ}へ2、そこから\R{上}{うえ}へ4\R{進}{すす}んだ\R{所}{ところ}にかきます（\R{青}{あお}い\R{点線}{てん|せん}）。ほかの\R{組}{くみ}も\R{同}{おな}じようにかきます。くわしくは3-2で\R{学}{まな}びます。

\textbf{まっすぐになるわけ}\quad $y=2x$ では、$x$ が1ふえるごとに、$y$ はいつも\R{同}{おな}じ2だけふえます（\R{左}{ひだり}のグラフの\R{緑}{みどり}の\R{階段}{かい|だん}）。\R{同}{おな}じ\R{形}{かたち}の\R{階段}{かい|だん}がくり\R{返}{かえ}すので、\R{点}{てん}はまっすぐにならびます。おもりが0\R{個}{こ}なら、のびも0なので、\R{線}{せん}は $x=0$、$y=0$ の\R{点}{てん}から\R{始}{はじ}まります。この\R{点}{てん}を\textbf{\R{原点}{げん|てん}}といいます。

\R{右}{みぎ}のグラフのように、はかった\R{点}{てん}（\R{黒}{くろ}）は、\R{原点}{げん|てん}から\R{始}{はじ}まるまっすぐな\R{線}{せん}（\R{青}{あお}）の\R{近}{ちか}くにならびます。\R{比例}{ひ|れい}のグラフは、\R{原点}{げん|てん}を\R{通}{とお}る\R{直線}{ちょく|せん}です。\R{比例定数}{ひ|れい|てい|すう}が\R{大}{おお}きいほど、1ふえるごとに\R{上}{あ}がる\R{高}{たか}さが\R{大}{おお}きいので、\R{線}{せん}は\R{急}{きゅう}になります（オレンジは $y=x$）。

\bksubsection{3}{ばねの\R{力}{ちから}：$F=kx$}

\leavevmode\label{bk:2-15:bane}\noindent\begin{minipage}[t]{0.47\linewidth}\vspace{0pt}
\R{物理}{ぶつ|り}では、おもりの\R{数}{かず}ではなく、ばねを\R{引}{ひ}く\R{力}{ちから}で\R{考}{かんが}えます。\R{力}{ちから}の\R{大}{おお}きさは、\textbf{N}（ニュートン）という\R{単位}{たん|い}ではかります（2-10）。
\end{minipage}\hfill
\begin{minipage}[t]{67mm}\vspace{0pt}\bkkigocard{\sffamily N}{ニュートン}{\R{力}{ちから}の\R{単位}{たん|い}。\R{質量}{しつ|りょう}\R{約}{やく}\qty{100}{g}の\R{物}{もの}にはたらく\R{重}{おも}さ（\R{重力}{じゅう|りょく}）が\R{約}{やく}\qty{1}{N}。}{newton}{2-15}\end{minipage}
\par\vspace{4pt}
\textbf{\qty{100}{g}で\R{約}{やく}\qty{1}{N}になるわけ}\quad \R{地球}{ち|きゅう}は、\R{質量}{しつ|りょう}\qty{1}{kg}の\R{物}{もの}を\R{約}{やく}\qty{10}{N}の\R{力}{ちから}で\R{下}{した}に\R{引}{ひ}いています（くわしくは\qty{9.8}{N}。2-0の「1\R{秒}{びょう}ごとに\R{約}{やく}\qty{10}{m/s}ずつ\R{速}{はや}くなる」の10と\R{同}{おな}じ\R{数}{かず}です。5-5）。この\R{力}{ちから}が、ふだん「\R{重}{おも}さ」とよんでいるものです。\qty{100}{g}は\qty{1}{kg}（\qty{1000}{g}）の10\R{分}{ぶん}の1なので、\R{引}{ひ}く\R{力}{ちから}も10\R{分}{ぶん}の1の\R{約}{やく}\qty{1}{N}です。\qty{50}{g}のおもり1\R{個}{こ}なら、その\R{半分}{はん|ぶん}の\R{約}{やく}\qty{0.5}{N}です。

\R{質量}{しつ|りょう}（g、kg）は「\R{物}{もの}そのものの\R{量}{りょう}」で、\R{月}{つき}に\R{持}{も}っていっても\R{変}{か}わりません。\R{重}{おも}さ（N）は「\R{地球}{ち|きゅう}が\R{引}{ひ}く\R{力}{ちから}」で、\R{月}{つき}では\R{約}{やく}6\R{分}{ぶん}の1になります。ばねがのびるのは、この\R{力}{ちから}で\R{引}{ひ}かれるからです。

\R{力}{ちから}を $F$\,N、のびを $x$\,m として、\R{表}{ひょう}を\R{作}{つく}り\R{直}{なお}します。のびはmにそろえます。\R{力}{ちから}の\R{単位}{たん|い}もmも、\R{世界}{せ|かい}\R{共通}{きょう|つう}のSI\R{単位}{たん|い}（2-10）だからです。cmをmに\R{直}{なお}すには、\qty{100}{cm/m}でわります。
\bkiipline{$\qty{2}{cm}\div\qty{100}{cm/m}=2\,\cancel{\mathrm{cm}}\times\dfrac{1\,\mathrm{m}}{100\,\cancel{\mathrm{cm}}}=\qty{0.02}{m}$}

\begin{bkcenter}
\begin{bktable}{colspec={Q[l,m,wd=36mm] *{4}{Q[c,m,wd=14mm]}}, rowsep=0.5pt}
おもりの\R{数}{かず}（\R{個}{こ}） & 1 & 2 & 3 & 4 \\
のび $x$（m） & 0.02 & 0.04 & 0.06 & 0.08 \\
\R{力}{ちから} $F$（N） & 0.5 & 1.0 & 1.5 & 2.0 \\
\SetCell{bg=bkYellow!30} $F\div x$（N/m） & 25 & 25 & 25 & 25 \\
\end{bktable}
\bktabcaption{\R{力}{ちから}は、のびに\R{比例}{ひ|れい}する。}
\end{bkcenter}

たとえば1\R{個}{こ}のとき、$\qty{0.5}{N}\div\qty{0.02}{m}=\qty{25}{N/m}$ です。$F\div x$ がいつも\R{同}{おな}じなので、\R{力}{ちから}はのびに\R{比例}{ひ|れい}します。\R{比例定数}{ひ|れい|てい|すう}を $k$ と\R{書}{か}くと、
\bkiipline{\fbox{$F=kx$}}
\noindent\begin{minipage}[c]{0.58\linewidth}
この $k$ を\textbf{ばね\R{定数}{てい|すう}}といいます。\bkidx{ばねていすう}{ばね\R{定数}{てい|すう}}\R{単位}{たん|い}は N/m（ニュートン\R{毎}{まい}メートル）で、「\qty{1}{m}あたり\R{何}{なん}Nいるか」という1あたりの\R{量}{りょう}です（\R{速}{はや}さ m/s が「1\R{秒}{びょう}あたり\R{何}{なん}m」なのと\R{同}{おな}じ）。ほんとうに\qty{1}{m}ものばすわけではありません。
\end{minipage}\hfill
\begin{minipage}[c]{0.39\linewidth}\centering
\begin{tikzpicture}
\begin{axis}[bkplot, width=44mm, height=46mm, xmin=0, xmax=0.095, ymin=0, ymax=2.4, xtick={0,0.04,0.08}, ytick={0,0.5,...,2.0},
  xticklabel style={/pgf/number format/fixed}, yticklabel style={/pgf/number format/fixed, /pgf/number format/precision=1},
  xlabel={$x$（m）}, ylabel={$F$（N）}]
  \addplot[bkcurve, domain=0:0.09] {25*x};
  \addplot[only marks, mark=*, mark size=1.7pt, bkInk] coordinates {(0.02,0.5) (0.04,1.0) (0.06,1.5) (0.08,2.0)};
  \node[font=\sffamily\footnotesize, text=bkBlue!80!black, anchor=east] at (axis cs:0.06,1.95) {$F=25x$};
\end{axis}
\end{tikzpicture}
\end{minipage}

\textbf{cmで\R{考}{かんが}えても\R{同}{おな}じ}\quad このばねを\qty{1}{cm}あたりで\R{考}{かんが}えると、
\[ \qty{0.5}{N}\div\qty{2}{cm}=\qty{0.25}{N/cm} \]
です。1 mの\R{中}{なか}には\qty{100}{cm}あるので、\qty{100}{cm/m}をかけると、cmが\R{約分}{やく|ぶん}されて\R{消}{き}えます。
\[ \qty{0.25}{N/cm}\times\qty{100}{cm/m}=\qty{25}{N/m} \]
2\R{通}{とお}りの\R{計算}{けい|さん}が\R{合}{あ}います。

\R{硬}{かた}いばねほど、\R{同}{おな}じだけのばすのに\R{大}{おお}きな\R{力}{ちから}がいるので、$k$ は\R{大}{おお}きくなります。\bkhatenaref のもう1つの\R{答}{こた}え：$F=kx$ で $x$ を2\R{倍}{ばい}にすると $F$ も2\R{倍}{ばい}なので、ばねを2\R{倍}{ばい}のばすには、2\R{倍}{ばい}の\R{力}{ちから}がいります。

\bksubsection{4}{\R{電圧}{でん|あつ}と\R{電流}{でん|りゅう}：オームの\R{法則}{ほう|そく}}

\leavevmode\label{bk:2-15:ohm}\noindent\begin{minipage}[c]{0.56\linewidth}
\R{抵抗}{てい|こう}（\R{電流}{でん|りゅう}を\R{流}{なが}れにくくする\R{部品}{ぶ|ひん}）に\R{乾電池}{かん|でん|ち}をつなぎ、\R{電流計}{でん|りゅう|けい}で\R{電流}{でん|りゅう}をはかりました。\R{電圧}{でん|あつ}（\R{電流}{でん|りゅう}をおし\R{出}{だ}すはたらきの\R{大}{おお}きさ）の\R{単位}{たん|い}は V（ボルト）、\R{電流}{でん|りゅう}の\R{単位}{たん|い}は A（アンペア）です。\R{乾電池}{かん|でん|ち}1\R{個}{こ}は\qty{1.5}{V}で、2\R{個}{こ}、3\R{個}{こ}と\R{直列}{ちょく|れつ}につなぐと（\R{小}{しょう}4の\R{理科}{り|か}）、\R{電圧}{でん|あつ}は $1.5\times2=3.0$、$1.5\times3=4.5$（V）とふえます。
\end{minipage}\hfill
\begin{minipage}[c]{0.4\linewidth}\centering
\begin{tikzpicture}[x=1mm, y=1mm]
  \draw[bkInk, line width=0.9pt] (0,9) -- (0,24) -- (40,24) -- (40,16.2) (40,7.8) -- (40,0) -- (0,0) -- (0,5.5);
  \draw[bkInk, line width=1.1pt] (-4.5,8.5) -- (4.5,8.5);
  \draw[bkInk, line width=2.2pt] (-2.2,5.5) -- (2.2,5.5);
  \node[font=\sffamily\scriptsize, text=bkInk, anchor=east] at (-4.8,10.3) {$+$};
  \filldraw[fill=bkYellowBg, draw=bkInk, line width=0.9pt] (12,21.6) rectangle (24,26.4);
  \filldraw[fill=white, draw=bkInk, line width=0.9pt] (40,12) circle (4.2);
  \node[font=\sffamily\bfseries\small, text=bkInk] at (40,12) {A};
  \draw[-{Stealth[length=5pt]}, bkBlue, line width=1.1pt] (27,27.5) -- (33,27.5);
  \node[font=\sffamily\footnotesize, text=bkBlue!80!black, anchor=south] at (30,29) {\R{電流}{でん|りゅう}};
  \node[font=\sffamily\footnotesize, text=bkInk, anchor=north] at (18,20) {\R{抵抗}{てい|こう}};
  \node[font=\sffamily\footnotesize, text=bkInk, anchor=west] at (5.5,7) {\R{乾電池}{かん|でん|ち}};
  \node[font=\sffamily\footnotesize, text=bkInk, anchor=east] at (34.5,12) {\R{電流計}{でん|りゅう|けい}};
\end{tikzpicture}
\end{minipage}

\begin{bkcenter}
\begin{bktable}{colspec={Q[l,m,wd=36mm] *{4}{Q[c,m,wd=14mm]}}, rowsep=0.5pt}
\R{乾電池}{かん|でん|ち}の\R{数}{かず}（\R{個}{こ}） & 1 & 2 & 3 & 4 \\
\R{電圧}{でん|あつ} $V$（V） & 1.5 & 3.0 & 4.5 & 6.0 \\
\R{電流}{でん|りゅう} $I$（A） & 0.03 & 0.06 & 0.09 & 0.12 \\
\SetCell{bg=bkYellow!30} $V\div I$（Ω） & 50 & 50 & 50 & 50 \\
\end{bktable}
\bktabcaption{\R{電流}{でん|りゅう}は、\R{電圧}{でん|あつ}に\R{比例}{ひ|れい}する。}
\end{bkcenter}

\noindent\begin{minipage}[t]{0.47\linewidth}\vspace{0pt}
\R{電圧}{でん|あつ}が2\R{倍}{ばい}なら\R{電流}{でん|りゅう}も2\R{倍}{ばい}。$V\div I$ はいつも50で\R{一定}{いっ|てい}です。この\R{値}{あたい}を、その\R{部品}{ぶ|ひん}の\textbf{\R{抵抗}{てい|こう}} $R$ といいます。\R{単位}{たん|い}は Ω（オーム）で、V÷A＝Ωです。
\end{minipage}\hfill
\begin{minipage}[t]{67mm}\vspace{0pt}\bkkigocard{Ω}{オーム}{\R{抵抗}{てい|こう}の\R{単位}{たん|い}。\qty{1}{V}で\qty{1}{A}\R{流}{なが}れる\R{抵抗}{てい|こう}が\qty{1}{\ohm}。}{ohm}{2-15}\end{minipage}
\par\vspace{4pt}
$R=\dfrac{V}{I}$ は、「\qty{1}{A}\R{流}{なが}すのに\R{何}{なん}Vいるか」という1あたりの\R{量}{りょう}です。\R{電流}{でん|りゅう}を\R{流}{なが}しにくい\R{部品}{ぶ|ひん}ほど、たくさんの\R{電圧}{でん|あつ}がいるので、$R$ は\R{大}{おお}きくなります。$R=\dfrac{V}{I}$ の\R{両辺}{りょう|へん}に $I$ をかけると $V=RI$、さらに\R{両辺}{りょう|へん}を $R$ でわると $I=\dfrac{V}{R}$ です（2-14）。この\R{関係}{かん|けい}を\textbf{オームの\R{法則}{ほう|そく}}といいます。\bkidx{おーむのほうそく}{オームの\R{法則}{ほう|そく}}$V=RI$ は、$F=kx$ と\R{同}{おな}じ\R{比例}{ひ|れい}の\R{形}{かたち}（$y=ax$）です。

\bksubsection{5}{かけると\R{一定}{いっ|てい}：\R{反比例}{はん|ぴ|れい}}

\leavevmode\label{bk:2-15:hanpirei}こんどは、\R{増}{ふ}えると\R{減}{へ}る2つの\R{量}{りょう}です。\qty{1200}{m}の\R{道}{みち}を\R{行}{い}くときの、\R{速}{はや}さ $x$（m/s）と、かかる\R{時間}{じ|かん} $y$（s）を\R{表}{ひょう}にします。かかる\R{時間}{じ|かん}は、\R{時間}{じ|かん}＝\R{道}{みち}のり÷\R{速}{はや}さ（2-14）で\R{計算}{けい|さん}します。たとえば $x=4$ なら $1200\div4=300$ です。

\begin{bkfigure}
\begin{tikzpicture}[x=1mm, y=1mm]
  \bkiiptable{\R{速}{はや}さ $x$（m/s）}{1,2,3,4,5,6}{\R{時間}{じ|かん} $y$（s）}{1200,600,400,300,240,200}
  \bkiiparrow{1}{0}{1}{2\R{倍}{ばい}}{bkBlue}
  \bkiiparrow{1}{1}{3}{2\R{倍}{ばい}}{bkBlue}
  \bkiiparrow{-1}{0}{1}{$\frac{1}{2}$\R{倍}{ばい}}{bkOrange}
  \bkiiparrow{-1}{1}{3}{$\frac{1}{2}$\R{倍}{ばい}}{bkOrange}
\end{tikzpicture}
\bkfigcaption{$x$ が2\R{倍}{ばい}になると、$y$ は $\dfrac{1}{2}$ \R{倍}{ばい}になる。}
\end{bkfigure}

\R{速}{はや}さが1から2へ2\R{倍}{ばい}になると、\R{時間}{じ|かん}は1200から600へ $\dfrac{1}{2}$ \R{倍}{ばい}。\R{速}{はや}さが1から3へ3\R{倍}{ばい}になると、\R{時間}{じ|かん}は1200から400へ $\dfrac{1}{3}$ \R{倍}{ばい}です。このように、$x$ が2\R{倍}{ばい}、3\R{倍}{ばい}、…になると、$y$ は $\dfrac{1}{2}$ \R{倍}{ばい}、$\dfrac{1}{3}$ \R{倍}{ばい}、…になるとき、「$y$ は $x$ に\textbf{\R{反比例}{はん|ぴ|れい}}する」といいます。\bkidx{はんぴれい}{\R{反比例}{はん|ぴ|れい}}

\noindent\begin{minipage}[c]{0.55\linewidth}
\textbf{かけ\R{算}{ざん}が\R{一定}{いっ|てい}}\quad \R{表}{ひょう}のたての\R{組}{くみ}で $x\times y$ を\R{計算}{けい|さん}すると、$1\times1200$、$2\times600$、$3\times400$、…で、どれも1200です。\R{速}{はや}さ×\R{時間}{じ|かん}は\R{道}{みち}のり（2-3）で、\R{道}{みち}のりはいつも\qty{1200}{m}だからです。
\end{minipage}\hfill
\begin{minipage}[c]{0.42\linewidth}\centering
\begin{tikzpicture}
\begin{axis}[bkplot, width=46mm, height=46mm, xmin=0, xmax=6.6, ymin=0, ymax=1400, xtick={0,...,6}, ytick={0,400,800,1200},
  xlabel={$x$（m/s）}, ylabel={$y$（s）}]
  \addplot[bkcurve, domain=0.88:6.4] {1200/x};
  \addplot[only marks, mark=*, mark size=1.6pt, bkInk] coordinates {(1,1200) (2,600) (3,400) (4,300) (5,240) (6,200)};
\end{axis}
\end{tikzpicture}
\end{minipage}

\vspace{3pt}
\R{一定}{いっ|てい}の\R{値}{あたい}を $a$ とすると、$xy=a$。\R{両辺}{りょう|へん}を $x$ でわって（2-14）、$y=\dfrac{a}{x}$ とも\R{書}{か}けます。$x$ を2\R{倍}{ばい}にすると、$\dfrac{a}{2x}$ で、\R{分母}{ぶん|ぼ}が2\R{倍}{ばい}になるので、$y$ は $\dfrac{1}{2}$ \R{倍}{ばい}になります。

\textbf{\R{反比例}{はん|ぴ|れい}のグラフ}\quad \R{点}{てん}をかきこむと、\R{上}{うえ}のグラフのように、\R{右}{みぎ}へ\R{行}{い}くほど\R{下}{さ}がり、\R{下}{さ}がり\R{方}{かた}がだんだんゆるやかになる\R{曲線}{きょく|せん}になります。\R{速}{はや}さをどれだけ\R{大}{おお}きくしても、\R{時間}{じ|かん}は0にはなりません。だから\R{曲線}{きょく|せん}は、\R{横}{よこ}の\R{軸}{じく}に\R{近}{ちか}づくけれど、くっつきません。\R{比例}{ひ|れい}とちがい、\R{直線}{ちょく|せん}ではなく、\R{原点}{げん|てん}も\R{通}{とお}りません（くわしくは3-7）。

\bksubsection{6}{2\R{倍}{ばい}にすると、どうなる？}

\leavevmode\label{bk:2-15:bai}「\R{増}{ふ}えると\R{増}{ふ}える」だけでは、\R{何倍}{なん|ばい}になるかはわかりません。$x$ を2\R{倍}{ばい}にしたとき $y$ が\R{何倍}{なん|ばい}になるかは、\R{関係}{かん|けい}の\R{形}{かたち}で\R{決}{き}まります。\R{比例}{ひ|れい}・\R{反比例}{はん|ぴ|れい}のほかに、これまでの\R{章}{しょう}で\R{出}{で}てきた2つの\R{関係}{かん|けい}もくらべてみましょう。

\textbf{2\R{乗}{じょう}に\R{比例}{ひ|れい}}\quad \R{物}{もの}が\R{落}{お}ちる\R{距離}{きょ|り}は $y=5x^2$（$x$ は\R{時間}{じ|かん}、2-13）。$x$ を2\R{倍}{ばい}にすると、$x^2$ は $2\times2=4$ \R{倍}{ばい}になります。$x=1$ で $y=5$、$x=2$ で $y=5\times4=20$、$x=4$ で $y=5\times16=80$。\R{時間}{じ|かん}が2\R{倍}{ばい}で4\R{倍}{ばい}、4\R{倍}{ばい}で16\R{倍}{ばい}です。$y=ax^2$ の\R{形}{かたち}を「$y$ は $x$ の2\R{乗}{じょう}に\R{比例}{ひ|れい}する」といいます。

\textbf{$\sqrt{\phantom{x}}$ に\R{比例}{ひ|れい}}\quad \R{振}{ふ}り\R{子}{こ}の\R{周期}{しゅう|き}は、\R{糸}{いと}の\R{長}{なが}さの $\sqrt{\phantom{x}}$ に\R{比例}{ひ|れい}します（2-8）。\R{長}{なが}さを4\R{倍}{ばい}にすると、\R{周期}{しゅう|き}は $\sqrt{4}=2$ \R{倍}{ばい}。\R{長}{なが}さを2\R{倍}{ばい}にすると、\R{周期}{しゅう|き}は $\sqrt{2}$ \R{倍}{ばい}で、\R{約}{やく}1.41\R{倍}{ばい}です。\R{反対}{はん|たい}に、\R{周期}{しゅう|き}を2\R{倍}{ばい}にしたければ、\R{長}{なが}さは $2\times2=4$ \R{倍}{ばい}にします。

\begin{bkcenter}
\begin{bktable}{colspec={Q[l,m,wd=31mm] Q[l,m,wd=56mm] Q[c,m,wd=17mm] Q[c,m,wd=17mm]}, rowsep=0pt, cells={font=\small}, row{1}={font=\sffamily\bfseries\small}}
\R{関係}{かん|けい} & \R{例}{れい} & 2\R{倍}{ばい}にする & 4\R{倍}{ばい}にする \\
\R{比例}{ひ|れい} $y=ax$ & ばねののびと\R{力}{ちから} & 2\R{倍}{ばい} & 4\R{倍}{ばい} \\
\R{反比例}{はん|ぴ|れい} $y=\dfrac{a}{x}$ & \R{速}{はや}さとかかる\R{時間}{じ|かん} & $\dfrac{1}{2}$\R{倍}{ばい} & $\dfrac{1}{4}$\R{倍}{ばい} \\
2\R{乗}{じょう}に\R{比例}{ひ|れい} $y=ax^2$ & \R{落}{お}ちる\R{時間}{じ|かん}と\R{距離}{きょ|り}（2-13） & 4\R{倍}{ばい} & 16\R{倍}{ばい} \\
$\sqrt{x}$ に\R{比例}{ひ|れい} & \R{糸}{いと}の\R{長}{なが}さと\R{振}{ふ}り\R{子}{こ}の\R{周期}{しゅう|き}（2-8） & \R{約}{やく}1.41\R{倍}{ばい} & 2\R{倍}{ばい} \\
\end{bktable}
\bktabcaption{$x$ を2\R{倍}{ばい}・4\R{倍}{ばい}にしたときの $y$。}
\end{bkcenter}

\begin{motto}[2]
\textbf{\R{注射器}{ちゅう|しゃ|き}の\R{空気}{くう|き}も\R{反比例}{はん|ぴ|れい}}\par
\R{先}{さき}をふさいだ\R{注射器}{ちゅう|しゃ|き}の\R{空気}{くう|き}を、\R{体積}{たい|せき}が\R{半分}{はん|ぶん}になるまでおしこむと、\R{中}{なか}の\R{空気}{くう|き}がおし\R{返}{かえ}す\R{強}{つよ}さ（\R{圧力}{あつ|りょく}）は2\R{倍}{ばい}になります。\R{温度}{おん|ど}が\R{変}{か}わらなければ、（おし\R{返}{かえ}す\R{強}{つよ}さ）×（\R{体積}{たい|せき}）が\R{一定}{いっ|てい}です。\R{空気}{くう|き}\R{入}{い}れ（ポンプ）のつつの\R{中}{なか}の\R{空気}{くう|き}も、このきまりにしたがいます（\R{高校}{こう|こう}の\R{物理}{ぶつ|り}で「ボイルの\R{法則}{ほう|そく}」として\R{学}{まな}びます）。
\end{motto}

\begin{column}{てこのつり\R{合}{あ}いも\R{反比例}{はん|ぴ|れい}}
\noindent\begin{minipage}[c]{0.55\linewidth}
てこがつり\R{合}{あ}うのは、（おもりの\R{重}{おも}さ）×（\R{支点}{し|てん}からの\R{距離}{きょ|り}）が\R{左右}{さ|ゆう}で\R{等}{ひと}しいときです（\R{小}{しょう}6の\R{理科}{り|か}）。\R{右}{みぎ}の\R{図}{ず}では、\R{左}{ひだり}が $2\times2=4$、\R{右}{みぎ}が $1\times4=4$。\R{支点}{し|てん}から2\R{倍}{ばい}はなすと、つり\R{合}{あ}う\R{重}{おも}さは $\dfrac{1}{2}$ になります。
\end{minipage}\hfill
\begin{minipage}[c]{0.42\linewidth}\centering
\begin{tikzpicture}[x=1mm, y=1mm]
  \fill[bkGray!70] (-2.5,-9) -- (2.5,-9) -- (0,-1) -- cycle;
  \fill[bkInk!75, rounded corners=0.6mm] (-26,-1) rectangle (26,1);
  \foreach \i in {-4,...,4} \fill[white] (\i*6,0) circle (0.5);
  \draw[bkInk!60, line width=0.5pt] (-12,-1) -- (-12,-5) (24,-1) -- (24,-5);
  \bkiipweight{(-12,-5)}\bkiipweight{(-12,-11.5)}
  \bkiipweight{(24,-5)}
  \node[bk note, anchor=north] at (-12,-18.5) {2\R{個}{こ}};
  \node[bk note, anchor=north] at (24,-12) {1\R{個}{こ}};
  \draw[{Stealth[length=3pt]}-{Stealth[length=3pt]}, bkBlue, line width=0.6pt] (-12,3) -- (0,3);
  \draw[{Stealth[length=3pt]}-{Stealth[length=3pt]}, bkBlue, line width=0.6pt] (0,3) -- (24,3);
  \node[font=\sffamily\scriptsize, text=bkBlue!80!black, anchor=south] at (-6,3.6) {2};
  \node[font=\sffamily\scriptsize, text=bkBlue!80!black, anchor=south] at (12,3.6) {4};
\end{tikzpicture}
\end{minipage}
\end{column}

\begin{taisetsu}
$y=ax$（\R{比例}{ひ|れい}）：$x$ が2\R{倍}{ばい}なら $y$ も2\R{倍}{ばい}。$y\div x$ が\R{一定}{いっ|てい}（＝$a$、1あたりの\R{量}{りょう}）。グラフは\R{原点}{げん|てん}を\R{通}{とお}る\R{直線}{ちょく|せん}。\par\vspace{2pt}
$y=\dfrac{a}{x}$（\R{反比例}{はん|ぴ|れい}）：$x$ が2\R{倍}{ばい}なら $y$ は $\dfrac{1}{2}$ \R{倍}{ばい}。$x\times y$ が\R{一定}{いっ|てい}（＝$a$）。\par\vspace{2pt}
\R{比例}{ひ|れい}か\R{反比例}{はん|ぴ|れい}かは、$y\div x$ や $x\times y$ を\R{計算}{けい|さん}して\R{確}{たし}かめる。
\end{taisetsu}

\begin{kotoba}
\small
\bkword{\R{比例}{ひ|れい}・\R{比例定数}{ひ|れい|てい|すう}}{$y=ax$ の\R{関係}{かん|けい}。$a=\dfrac{y}{x}$ が\R{比例定数}{ひ|れい|てい|すう}。グラフは\R{原点}{げん|てん}を\R{通}{とお}る\R{直線}{ちょく|せん}。}
\bkword{\R{反比例}{はん|ぴ|れい}}{$y=\dfrac{a}{x}$ の\R{関係}{かん|けい}。$xy$ が\R{一定}{いっ|てい}。}
\bkword{ばね\R{定数}{てい|すう}}{$F=kx$ の $k$。\qty{1}{m}のばすのに\R{何}{なん}Nいるか（1あたりの\R{量}{りょう}）。}
\bkword{オームの\R{法則}{ほう|そく}}{$V=RI$。\R{電流}{でん|りゅう}は\R{電圧}{でん|あつ}に\R{比例}{ひ|れい}する。$R$ は\R{抵抗}{てい|こう}。}
\end{kotoba}

%% ---------------- 例題 ----------------
\bksection{reidai}{\R{例題}{れい|だい}でたしかめよう}

\begin{reidai}[\R{比例}{ひ|れい}の\R{式}{しき}を\R{作}{つく}る]
はり\R{金}{がね}の\R{重}{おも}さ $y$\,g は\R{長}{なが}さ $x$\,m に\R{比例}{ひ|れい}し、\qty{2}{m}で\qty{30}{g}です。$y$ を $x$ の\R{式}{しき}で\R{表}{あらわ}し、\qty{8}{m}の\R{重}{おも}さを\R{求}{もと}めましょう。\par
\kangaekata \R{比例}{ひ|れい}なので $y=ax$ の\R{形}{かたち}です。\R{比例定数}{ひ|れい|てい|すう} $a=\dfrac{y}{x}$ は、1\,mあたりの\R{重}{おも}さです。\par
\shiki $a=\dfrac{30}{2}=15$\quad わかっている $x=2$、$y=30$ を\R{使}{つか}う\par
\hspace*{2.6\zw}だから $y=15x$。$x=8$ を\R{代入}{だい|にゅう}して $y=15\times8=120$\par
\kotae $y=15x$、\qty{120}{g}（\R{確}{たし}かめ：\R{長}{なが}さが4\R{倍}{ばい}なので\R{重}{おも}さも4\R{倍}{ばい}、$30\times4=120$）
\end{reidai}

\begin{reidai}[ばねの\R{力}{ちから} $F=kx$]
（1）\qty{1.5}{N}で\R{引}{ひ}くと\qty{5}{cm}のびるばねの、ばね\R{定数}{てい|すう}は\R{何}{なん}N/mですか。\par
（2）ばね\R{定数}{てい|すう}\qty{40}{N/m}のばねを、\qty{5}{cm}のばす\R{力}{ちから}は\R{何}{なん}Nですか。\par
（3）（2）のばねを\qty{3}{N}で\R{引}{ひ}くと、\R{何}{なん}cmのびますか。\par
\kangaekata ばね\R{定数}{てい|すう}は「\qty{1}{m}あたり」なので、のびはmにそろえます。（1）は $k=\dfrac{F}{x}$、（3）は $x=\dfrac{F}{k}$（$F=kx$ を2-14のように\R{変形}{へん|けい}）を\R{使}{つか}います。\par
\shiki $\qty{5}{cm}\div\qty{100}{cm/m}=\qty{0.05}{m}$\par
\hspace*{2.6\zw}（1）\nobreak$k=\dfrac{1.5}{0.05}=30$\quad（2）\nobreak$F=40\times0.05=2$\par
\hspace*{2.6\zw}（3）\nobreak$x=\dfrac{3}{40}=0.075$。$\qty{0.075}{m}\times\qty{100}{cm/m}=\qty{7.5}{cm}$\par
\kotae（1）\nobreak\qty{30}{N/m}\quad（2）\nobreak\qty{2}{N}\quad（3）\nobreak\qty{7.5}{cm}
\end{reidai}

\begingroup\tcbset{equal height group=bkiipjA}%
\noindent\begin{minipage}[t]{0.49\linewidth}
\begin{jibun}
\R{水}{みず}の\R{深}{ふか}さ $y$\,cm は\R{時間}{じ|かん} $x$ \R{分}{ふん}に\R{比例}{ひ|れい}し、3\R{分}{ぷん}で\qty{12}{cm}でした。$y$ を $x$ の\R{式}{しき}で\R{表}{あらわ}し、10\R{分後}{ぷん|ご}の\R{深}{ふか}さを\R{求}{もと}めましょう。
\bkjans{$a=\dfrac{12}{3}=4$ で $y=4x$。$y=4\times10=40$ で\qty{40}{cm}}
\end{jibun}
\end{minipage}\hfill
\begin{minipage}[t]{0.49\linewidth}
\begin{jibun}
（1）\qty{1.5}{N}で\qty{6}{cm}のびるばねの、ばね\R{定数}{てい|すう}を\R{求}{もと}めましょう。（2）このばねを\qty{12}{cm}のばす\R{力}{ちから}と、（3）\qty{2}{N}で\R{引}{ひ}いたときののびを\R{求}{もと}めましょう。
\bkjans{（1）\nobreak $\qty{6}{cm}\div\qty{100}{cm/m}=\qty{0.06}{m}$、$k=\dfrac{1.5}{0.06}=25$ で\qty{25}{N/m}\quad（2）\nobreak $25\times0.12=3$ で\qty{3}{N}\quad（3）\nobreak $\dfrac{2}{25}=0.08$ で\qty{0.08}{m}、$\qty{0.08}{m}\times\qty{100}{cm/m}=\qty{8}{cm}$}
\end{jibun}
\end{minipage}\par
\endgroup

\begin{reidai}[オームの\R{法則}{ほう|そく}]
（1）\qty{20}{\ohm}の\R{抵抗}{てい|こう}に\qty{6}{V}をかけると\R{何}{なん}A\R{流}{なが}れますか。（2）\qty{20}{\ohm}の\R{抵抗}{てい|こう}に\qty{0.5}{A}\R{流}{なが}すには\R{何}{なん}Vかければよいですか。（3）\qty{6}{V}をかけると\qty{0.3}{A}\R{流}{なが}れる\R{部品}{ぶ|ひん}の\R{抵抗}{てい|こう}は\R{何}{なん}Ωですか。\par
\kangaekata ほしい\R{量}{りょう}について\R{解}{と}いた\R{式}{しき}を\R{使}{つか}います。（1）は $I=\dfrac{V}{R}$、（2）は $V=RI$、（3）は $R=\dfrac{V}{I}$。\par
\shiki（1）\nobreak$I=\dfrac{6}{20}=0.3$\quad（2）\nobreak$V=20\times0.5=10$\quad（3）\nobreak$R=\dfrac{6}{0.3}=20$\par
\kotae（1）\nobreak\qty{0.3}{A}\quad（2）\nobreak\qty{10}{V}\quad（3）\nobreak\qty{20}{\ohm}
\end{reidai}

\begin{reidai}[\R{反比例}{はん|ぴ|れい}と、2\R{倍}{ばい}・$\dfrac{1}{2}$ \R{倍}{ばい}にすると？]
\qty{4}{m/s}で300\R{秒}{びょう}かかる\R{道}{みち}を、（1）\qty{6}{m/s}で\R{行}{い}くと\R{何秒}{なん|びょう}かかりますか。（2）\R{速}{はや}さを2\R{倍}{ばい}にすると、\R{時間}{じ|かん}は\R{何倍}{なん|ばい}ですか。（3）\R{速}{はや}さを $\dfrac{1}{2}$ \R{倍}{ばい}にすると、\R{時間}{じ|かん}は\R{何倍}{なん|ばい}ですか。\par
\kangaekata \R{道}{みち}のりが\R{同}{おな}じなので、\R{速}{はや}さ×\R{時間}{じ|かん}が\R{一定}{いっ|てい}（\R{反比例}{はん|ぴ|れい}）です。\par
\shiki $xy=4\times300=1200$\quad（1）\nobreak$y=\dfrac{1200}{6}=200$\par
\hspace*{2.6\zw}（2）\nobreak\qty{8}{m/s}なら $y=\dfrac{1200}{8}=150$。$150\div300=\dfrac{1}{2}$\par
\hspace*{2.6\zw}（3）\nobreak\qty{2}{m/s}なら $y=\dfrac{1200}{2}=600$。$600\div300=2$\par
\kotae（1）\nobreak200\R{秒}{びょう}\quad（2）\nobreak$\dfrac{1}{2}$\R{倍}{ばい}\quad（3）\nobreak2\R{倍}{ばい}
\end{reidai}

\begingroup\tcbset{equal height group=bkiipjB}%
\noindent\begin{minipage}[t]{0.49\linewidth}
\begin{jibun}
（1）\qty{50}{\ohm}の\R{抵抗}{てい|こう}に\qty{4.5}{V}をかけたときの\R{電流}{でん|りゅう}、（2）\qty{50}{\ohm}に\qty{0.2}{A}\R{流}{なが}すのに\R{必要}{ひつ|よう}な\R{電圧}{でん|あつ}、（3）\qty{9}{V}で\qty{0.3}{A}\R{流}{なが}れる\R{部品}{ぶ|ひん}の\R{抵抗}{てい|こう}を\R{求}{もと}めましょう。
\bkjans{（1）\nobreak $I=\dfrac{4.5}{50}=0.09$ で\qty{0.09}{A}\quad（2）\nobreak $V=50\times0.2=10$ で\qty{10}{V}\quad（3）\nobreak $R=\dfrac{9}{0.3}=30$ で\qty{30}{\ohm}}
\end{jibun}
\end{minipage}\hfill
\begin{minipage}[t]{0.49\linewidth}
\begin{jibun}
かみ\R{合}{あ}う\R{歯車}{は|ぐるま}は、（\R{歯}{は}の\R{数}{かず}）×（\R{回転数}{かい|てん|すう}）が\R{等}{ひと}しくなります（かみ\R{合}{あ}った\R{所}{ところ}を\R{通}{とお}る\R{歯}{は}の\R{数}{かず}が\R{同}{おな}じだから）。（1）\R{歯}{は}が30の\R{歯車}{は|ぐるま}が4\R{回転}{かい|てん}すると、\R{歯}{は}が20の\R{歯車}{は|ぐるま}は\R{何回転}{なん|かい|てん}しますか。（2）\R{歯}{は}の\R{数}{かず}を $\dfrac{1}{2}$ \R{倍}{ばい}にすると、\R{回転数}{かい|てん|すう}は\R{何倍}{なん|ばい}ですか。
\bkjans{（1）\nobreak $30\times4=120$、$120\div20=6$ で6\R{回転}{かい|てん}\quad（2）\nobreak 2\R{倍}{ばい}}
\end{jibun}
\end{minipage}\par
\endgroup

\begin{reidai}[\R{比例}{ひ|れい}・\R{反比例}{はん|ぴ|れい}・どちらでもない、を\R{見分}{み|わ}ける]
$y$ を $x$ の\R{式}{しき}で\R{表}{あらわ}し、\R{比例}{ひ|れい}・\R{反比例}{はん|ぴ|れい}・どちらでもない、のどれかを\R{答}{こた}えましょう。\par
（1）1 mが120\R{円}{えん}のリボンを $x$\,m\R{買}{か}ったときの\R{代金}{だい|きん} $y$ \R{円}{えん}\par
（2）\R{面積}{めん|せき}が\qty{18}{cm^2}の\R{長方形}{ちょう|ほう|けい}の、たて $x$\,cm と\R{横}{よこ} $y$\,cm\par
（3）10\,Lの\R{水}{みず}が\R{入}{はい}った\R{水}{みず}そうに、1\R{分}{ぷん}に2\,Lずつ\R{水}{みず}を\R{入}{い}れたときの、$x$ \R{分後}{ふん|ご}の\R{水}{みず}の\R{量}{りょう} $y$\,L\par
\kangaekata \R{式}{しき}を\R{作}{つく}ってから、$x=1$、2、3 を\R{代入}{だい|にゅう}して\R{表}{ひょう}にし、$y\div x$ と $x\times y$ を\R{計算}{けい|さん}します。\par
\shiki（1）\nobreak $y=120x$。$x=1,2,3$ で $y=120,240,360$。$y\div x$ はいつも120\par
\hspace*{2.6\zw}（2）\nobreak $xy=18$ より $y=\dfrac{18}{x}$。$x=1,2,3$ で $y=18,9,6$。$x\times y$ はいつも18\par
\hspace*{2.6\zw}（3）\nobreak $y=2x+10$。$x=1,2,3$ で $y=12,14,16$。$y\div x$ は12、7、5.33…でばらばら。$x\times y$ も12、28、48でばらばら\par
\kotae（1）\nobreak \R{比例}{ひ|れい}\quad（2）\nobreak \R{反比例}{はん|ぴ|れい}\quad（3）\nobreak どちらでもない（ふえるけれど、2\R{倍}{ばい}にはならない）
\end{reidai}

\begin{reidai}[2\R{乗}{じょう}に\R{比例}{ひ|れい}・$\sqrt{\phantom{x}}$ に\R{比例}{ひ|れい}]
（1）\R{物}{もの}が\R{落}{お}ちる\R{距離}{きょ|り}は、\R{時間}{じ|かん}の2\R{乗}{じょう}に\R{比例}{ひ|れい}します。1\R{秒}{びょう}で\qty{5}{m}\R{落}{お}ちるとき、3\R{秒}{びょう}では\R{何倍}{なん|ばい}、\R{何}{なん}m\R{落}{お}ちますか。\par
（2）\R{振}{ふ}り\R{子}{こ}の\R{周期}{しゅう|き}は、\R{糸}{いと}の\R{長}{なが}さの $\sqrt{\phantom{x}}$ に\R{比例}{ひ|れい}します。\R{糸}{いと}の\R{長}{なが}さ\R{約}{やく}\qty{25}{cm}で\R{周期}{しゅう|き}が1\R{秒}{びょう}の\R{振}{ふ}り\R{子}{こ}を、\R{周期}{しゅう|き}3\R{秒}{びょう}にするには、\R{長}{なが}さを\R{何倍}{なん|ばい}、\R{約}{やく}\R{何}{なん}cmにしますか。\par
\kangaekata （1）\R{時間}{じ|かん}が3\R{倍}{ばい}なら、2\R{乗}{じょう}は $3\times3=9$ \R{倍}{ばい}。（2）\R{周期}{しゅう|き}を3\R{倍}{ばい}にするには、$\sqrt{\phantom{x}}$ が3\R{倍}{ばい}になる\R{長}{なが}さ、つまり $3\times3=9$ \R{倍}{ばい}の\R{長}{なが}さにします（$\sqrt{9}=3$。2-8）。\par
\shiki（1）\nobreak $5\times9=45$\quad（2）\nobreak $25\times9=225$\par
\kotae（1）\nobreak 9\R{倍}{ばい}、\qty{45}{m}\quad（2）\nobreak 9\R{倍}{ばい}、\R{約}{やく}\qty{225}{cm}
\end{reidai}

\begingroup\tcbset{equal height group=bkiipjC}%
\noindent\begin{minipage}[t]{0.49\linewidth}
\begin{jibun}
$y$ を $x$ の\R{式}{しき}で\R{表}{あらわ}し、\R{比例}{ひ|れい}・\R{反比例}{はん|ぴ|れい}・どちらでもない、のどれかを\R{答}{こた}えましょう。（1）1ふくろ6\R{個}{こ}\R{入}{い}りのあめを $x$ ふくろ\R{買}{か}ったときのあめの\R{数}{かず} $y$ \R{個}{こ}（2）30\R{個}{こ}のあめを $x$ \R{人}{にん}で\R{同}{おな}じ\R{数}{かず}ずつ\R{分}{わ}けたときの1\R{人分}{にん|ぶん} $y$ \R{個}{こ}（3）100ページの\R{本}{ほん}を $x$ ページ\R{読}{よ}んだときの\R{残}{のこ}り $y$ ページ
\bkjans{（1）\nobreak $y=6x$、\R{比例}{ひ|れい}\quad（2）\nobreak $y=\dfrac{30}{x}$、\R{反比例}{はん|ぴ|れい}\quad（3）\nobreak $y=100-x$、どちらでもない}
\end{jibun}
\end{minipage}\hfill
\begin{minipage}[t]{0.49\linewidth}
\begin{jibun}
（1）1\R{秒}{びょう}で\qty{5}{m}\R{落}{お}ちる\R{物}{もの}は、4\R{秒}{びょう}では\R{何倍}{なん|ばい}、\R{何}{なん}m\R{落}{お}ちますか。（2）\R{周期}{しゅう|き}1\R{秒}{びょう}の\R{振}{ふ}り\R{子}{こ}（\R{約}{やく}\qty{25}{cm}）の\R{糸}{いと}を16\R{倍}{ばい}の\R{約}{やく}\qty{400}{cm}にすると、\R{周期}{しゅう|き}は\R{何秒}{なん|びょう}になりますか。
\bkjans{（1）\nobreak $4\times4=16$ \R{倍}{ばい}、$5\times16=80$ で\qty{80}{m}\quad（2）\nobreak $\sqrt{16}=4$ で4\R{倍}{ばい}、\R{約}{やく}4\R{秒}{びょう}}
\end{jibun}
\end{minipage}\par
\endgroup

\begin{reidai}[\R{比例}{ひ|れい}を\R{使}{つか}って、\R{数}{かぞ}えずに\R{数}{かず}を\R{求}{もと}める]
\R{同}{おな}じ\R{紙}{かみ}100\R{枚}{まい}の\R{重}{おも}さをはかると\qty{400}{g}でした。この\R{紙}{かみ}のたばの\R{重}{おも}さが\qty{2.6}{kg}のとき、\R{紙}{かみ}は\R{何枚}{なん|まい}ありますか。\par
\kangaekata \R{同}{おな}じ\R{紙}{かみ}なら、\R{重}{おも}さ $y$\,g は\R{枚数}{まい|すう} $x$ \R{枚}{まい}に\R{比例}{ひ|れい}します。1\R{枚}{まい}あたりの\R{重}{おも}さ（\R{比例定数}{ひ|れい|てい|すう}）を\R{求}{もと}めます。\R{重}{おも}さはgにそろえます（\qty{1000}{g/kg}をかける）。\par
\shiki $a=\dfrac{400}{100}=4$ だから $y=4x$\quad 1\R{枚}{まい}\qty{4}{g}\par
\hspace*{2.6\zw}$\qty{2.6}{kg}\times\qty{1000}{g/kg}=\qty{2600}{g}$\quad kgが\R{約分}{やく|ぶん}されてg\par
\hspace*{2.6\zw}$2600=4x$ の\R{両辺}{りょう|へん}を4でわって（2-14）、$x=650$\par
\kotae 650\R{枚}{まい}（\R{確}{たし}かめ：$4\times650=2600$）
\end{reidai}

\begin{jibun}
\R{同}{おな}じくぎ50\R{本}{ほん}の\R{重}{おも}さは\qty{120}{g}でした。このくぎ\R{全部}{ぜん|ぶ}の\R{重}{おも}さが\qty{0.96}{kg}のとき、くぎは\R{何本}{なん|ぼん}ありますか。
\bkjans{1\R{本}{ぽん}は $\dfrac{120}{50}=2.4$ で\qty{2.4}{g}、$y=2.4x$。$\qty{0.96}{kg}\times\qty{1000}{g/kg}=\qty{960}{g}$。$x=\dfrac{960}{2.4}=400$ で400\R{本}{ぼん}}
\end{jibun}

%% ---------------- まちがえやすいところ・物理工学 ----------------
\bksection{tsumazuki}{まちがえやすいところをチェック}\label{bk:2-15:tsuma}

\begin{tsumazuki}{「\R{増}{ふ}えると\R{増}{ふ}える」なら\R{比例}{ひ|れい}？}
\bkng{ばねの\R{全体}{ぜん|たい}の\R{長}{なが}さは、おもりの\R{数}{かず}に\R{比例}{ひ|れい}する}
\bkok{\R{比例}{ひ|れい}するのは「のび」。\R{全体}{ぜん|たい}の\R{長}{なが}さは\R{比例}{ひ|れい}しない}
もとの\R{長}{なが}さ\qty{12}{cm}なら、1\R{個}{こ}で14、2\R{個}{こ}で\qty{16}{cm}。2\R{倍}{ばい}になりません。$y\div x$ が\R{一定}{いっ|てい}かで\R{確}{たし}かめよう。
\end{tsumazuki}

\begin{tsumazuki}{グラフが\R{直線}{ちょく|せん}なら\R{比例}{ひ|れい}？}
\bkng{$y=x+2$ のグラフはまっすぐなので、\R{比例}{ひ|れい}}
\bkok{$x=0$ で $y=2$。\R{原点}{げん|てん}を\R{通}{とお}らないので、\R{比例}{ひ|れい}ではない}
\R{比例}{ひ|れい}のグラフは、\R{直線}{ちょく|せん}で、しかも\R{原点}{げん|てん}を\R{通}{とお}ります。$x=1$ で3、$x=2$ で4。2\R{倍}{ばい}になっていません。
\end{tsumazuki}

\begin{tsumazuki}{ばね\R{定数}{てい|すう}の\R{単位}{たん|い}をそろえない}
\bkng{\qty{0.5}{N}で\qty{2}{cm}のびたので、$k=0.5\div2=0.25$\,N/m}
\bkok{$\qty{2}{cm}\div\qty{100}{cm/m}=\qty{0.02}{m}$ だから、$k=0.5\div0.02=25$\,N/m}
N/m は「\qty{1}{m}あたり」です。cmのままだと100\R{倍}{ばい}ずれます（2-10）。
\end{tsumazuki}

\begin{tsumazuki}{「\R{増}{ふ}えると\R{減}{へ}る」なら\R{反比例}{はん|ぴ|れい}？}
\bkng{\qty{10}{L}から $x$\,L\R{使}{つか}った\R{残}{のこ}り $y=10-x$ は、\R{反比例}{はん|ぴ|れい}}
\bkok{$x=1$ で $y=9$、$x=2$ で $y=8$。$xy$ が\R{一定}{いっ|てい}でないので\R{反比例}{はん|ぴ|れい}でない}
\R{反比例}{はん|ぴ|れい}かどうかは、$x\times y$ をかけて\R{確}{たし}かめよう。
\end{tsumazuki}


\begin{butsuri}
\noindent\begin{minipage}[c]{0.72\linewidth}
\par\addvspace{3pt}\textbf{ばねばかり}\par\noindent ばねばかりの\R{目}{め}もりが\R{同}{おな}じ\R{間隔}{かん|かく}でならんでいるのは、のびが\R{力}{ちから}に\R{比例}{ひ|れい}するからです。\R{力}{ちから}が1 Nふえるごとに、のびはいつも\R{同}{おな}じだけふえます。のびをはかれば、$F=kx$ で\R{力}{ちから}がわかります。

\par\addvspace{3pt}\textbf{\R{自動車}{じ|どう|しゃ}のばね}\par\noindent タイヤの\R{近}{ちか}くには、\R{車}{くるま}の\R{重}{おも}さを\R{支}{ささ}えて、でこぼこ\R{道}{みち}のゆれをやわらげるばねがあります。ばね\R{定数}{てい|すう}が\R{約}{やく}\qty{30000}{N/m}もある、\R{硬}{かた}いばねです。1\R{本}{ぽん}のばねに、さらに\qty{150}{N}（\R{約}{やく}\qty{15}{kg}の\R{荷物}{に|もつ}の\R{重}{おも}さ）がかかると、
\bkiipline{$x=\dfrac{150}{30000}=0.005$、\quad $\qty{0.005}{m}\times\qty{1000}{mm/m}=\qty{5}{mm}$}
で、\qty{5}{mm}\R{沈}{しず}みます。
\end{minipage}\hfill
\begin{minipage}[c]{0.24\linewidth}\centering
\begin{tikzpicture}[x=1mm, y=1mm]
  \draw[bkInk, line width=0.8pt] (0,2) -- (0,6) (-2.5,6) arc (180:0:2.5);
  \filldraw[fill=bkBlue!10, draw=bkBlue!70!black, line width=0.8pt, rounded corners=1.5mm] (-6,-40) rectangle (6,2);
  \foreach \k in {0,...,5} {\draw[bkInk, line width=0.6pt] (-6,{-5-\k*6}) -- (-3,{-5-\k*6});
    \node[font=\sffamily\scriptsize, text=bkInk, anchor=west] at (0,{-5-\k*6}) {\k};}
  \draw[bkRed, line width=1.2pt] (-5.5,-23) -- (-1.5,-23);
  \draw[bkInk, line width=0.8pt] (0,-40) -- (0,-45);
  \draw[bkInk, line width=0.8pt] (-1.8,-45) arc (180:360:1.8);
  \node[bk note, anchor=west] at (7.5,-20) {N};
\end{tikzpicture}
\end{minipage}

\vspace{2pt}
\par\addvspace{3pt}\textbf{\R{自転車}{じ|てん|しゃ}のギア}\par\noindent \R{前}{まえ}のギアの\R{歯}{は}が44\R{枚}{まい}のとき、ペダル1\R{回転}{かい|てん}で、チェーンは\R{歯}{は}44\R{枚}{まい}\R{分}{ぶん}\R{進}{すす}みます。\R{後}{うし}ろのギアは、\R{歯}{は}が11\R{枚}{まい}なら $44\div11=4$ で4\R{回転}{かい|てん}、22\R{枚}{まい}なら $44\div22=2$ で2\R{回転}{かい|てん}します。（\R{歯}{は}の\R{数}{かず}）×（\R{回転数}{かい|てん|すう}）＝44で\R{一定}{いっ|てい}なので、\R{後}{うし}ろのギアの\R{回転数}{かい|てん|すう}は、\R{歯}{は}の\R{数}{かず}に\R{反比例}{はん|ぴ|れい}します。\R{歯}{は}の\R{少}{すく}ないギアほど\R{速}{はや}く\R{走}{はし}れますが、ペダルは\R{重}{おも}くなります。

\vspace{2pt}
\par\addvspace{3pt}\textbf{\R{大学}{だい|がく}への\R{種}{たね}}\par\noindent ばねの $F=kx$ と\R{電気}{でん|き}の $V=RI$ は、どちらも\R{比例}{ひ|れい}の\R{形}{かたち}です。ばねにつるしたおもりのゆれ（6-6）と、\R{電気}{でん|き}の\R{回路}{かい|ろ}の\R{中}{なか}のゆれも、\R{同}{おな}じ\R{形}{かたち}の\R{式}{しき}で\R{書}{か}けます。\R{機械}{き|かい}と\R{電気}{でん|き}を、1つの\R{数学}{すう|がく}で\R{設計}{せっ|けい}できるのです。
\end{butsuri}


%% ---------------- 章末問題 ----------------
\bkendproblems
\bkproblemgroup{\R{基本}{き|ほん}}{1}

\bkproblem $y$ を $x$ の\R{式}{しき}で\R{表}{あらわ}し、\R{比例}{ひ|れい}・\R{反比例}{はん|ぴ|れい}・どちらでもない、のどれかを\R{答}{こた}えましょう（$y\div x$ や $x\times y$ を\R{計算}{けい|さん}して\R{確}{たし}かめよう。\pageref{bk:2-15:tsuma}ページ）。\par
\hspace*{1\zw}（1）\nobreak 1\R{個}{こ}80\R{円}{えん}のパンを $x$ \R{個}{こ}\R{買}{か}ったときの\R{代金}{だい|きん} $y$ \R{円}{えん}\par
\hspace*{1\zw}（2）\nobreak \qty{24}{cm^2}の\R{長方形}{ちょう|ほう|けい}の、たて $x$\,cm と\R{横}{よこ} $y$\,cm\par
\hspace*{1\zw}（3）\nobreak 20\R{歳}{さい}の\R{兄}{あに}の、$x$ \R{年後}{ねん|ご}の\R{年齢}{ねん|れい} $y$ \R{歳}{さい}
\bkpans{（1）\nobreak $y=80x$、\R{比例}{ひ|れい}\quad（2）\nobreak $y=\dfrac{24}{x}$、\R{反比例}{はん|ぴ|れい}\quad（3）\nobreak $y=x+20$、どちらでもない}

\bkproblem $y$ は $x$ に\R{比例}{ひ|れい}し、$x=4$ のとき $y=10$ です。$y$ を $x$ の\R{式}{しき}で\R{表}{あらわ}し、$x=10$ のときの $y$ を\R{求}{もと}めましょう。
\bkpans{$a=\dfrac{10}{4}=2.5$ で $y=2.5x$。$y=2.5\times10=25$}

\bkproblem ばね\R{定数}{てい|すう}\qty{50}{N/m}のばねがあります。（1）\qty{4}{cm}のばす\R{力}{ちから}は\R{何}{なん}Nですか。（2）\qty{3}{N}で\R{引}{ひ}くと\R{何}{なん}cmのびますか。
\bkpans{（1）\nobreak $\qty{4}{cm}\div\qty{100}{cm/m}=\qty{0.04}{m}$、$50\times0.04=2$ で\qty{2}{N}\quad（2）\nobreak $\dfrac{3}{50}=0.06$ で\qty{0.06}{m}、$\qty{0.06}{m}\times\qty{100}{cm/m}=\qty{6}{cm}$}

\bkproblem \qty{10}{\ohm}の\R{抵抗}{てい|こう}があります。（1）\qty{3}{V}をかけたときの\R{電流}{でん|りゅう}は\R{何}{なん}Aですか。（2）\qty{1.2}{A}\R{流}{なが}すには\R{何}{なん}Vかければよいですか。
\bkpans{（1）\nobreak $\dfrac{3}{10}=0.3$ で\qty{0.3}{A}\quad（2）\nobreak $10\times1.2=12$ で\qty{12}{V}}

\bkproblem \qty{36}{km}はなれた\R{町}{まち}へ\R{行}{い}きます。（1）\R{時速}{じ|そく}\qty{12}{km}で\R{行}{い}くと\R{何時間}{なん|じ|かん}かかりますか。（2）2\R{時間}{じ|かん}で\R{着}{つ}くには、\R{時速}{じ|そく}\R{何}{なん}kmで\R{行}{い}けばよいですか。
\bkpans{（1）\nobreak $36\div12=3$ で3\R{時間}{じ|かん}\quad（2）\nobreak $36\div2=18$ で\R{時速}{じ|そく}\qty{18}{km}}
\begin{bkfigure}
\begin{tikzpicture}[x=1mm, y=1mm]
  \draw[bkGray!60, line width=3pt] (0,0) -- (108,0);
  \foreach \k in {0,...,6} {\draw[bkInk, line width=0.7pt] (\k*18,-1.6) -- (\k*18,1.6);
    \pgfmathtruncatemacro\km{\k*6}\node[bk tick, below=2mm] at (\k*18,0) {\km};}
  \node[bk tick, anchor=north west] at (110,-2.5) {km};
  % 家
  \fill[bkOrange!70] (-4,3) rectangle (4,9); \fill[bkRed!70] (-5.5,9) -- (0,13.5) -- (5.5,9) -- cycle;
  \node[bk label, anchor=south] at (0,14.5) {\R{家}{いえ}};
  % 町（ビル）
  \foreach \dx/\h in {-6/12,0/16,6/9} \fill[bkBlue!55] (108+\dx-2.5,3) rectangle (108+\dx+2.5,3+\h);
  \node[bk label, anchor=south] at (108,20) {\R{町}{まち}};
  \bkcar{(30,2)}{0.45}{bkGreen}
  \draw[bk arrow, bkGreen] (46,7) -- (58,7);
  \draw[{Stealth[length=4pt]}-{Stealth[length=4pt]}, bkRed, line width=0.7pt] (0,-10.5) -- (108,-10.5);
  \node[font=\sffamily\footnotesize, text=bkRed!80!black, fill=white, inner sep=1pt] at (54,-10.5) {\qty{36}{km}};
\end{tikzpicture}
\bkfigcaption{\R{速}{はや}さ×\R{時間}{じ|かん}は、いつも\qty{36}{km}。}
\end{bkfigure}

\bkproblem $y$ が $x$ の2\R{乗}{じょう}に\R{比例}{ひ|れい}するとき、$x$ を3\R{倍}{ばい}にすると $y$ は\R{何倍}{なん|ばい}になりますか。$y$ が $x$ に\R{反比例}{はん|ぴ|れい}するときはどうですか。
\bkpans{2\R{乗}{じょう}に\R{比例}{ひ|れい}：$3\times3=9$ で9\R{倍}{ばい}。\R{反比例}{はん|ぴ|れい}：$\dfrac{1}{3}$ \R{倍}{ばい}}

\bkproblemgroup{チャレンジ}{2}

\bkproblem \noindent\begin{minipage}[t]{0.52\linewidth}
\R{周期}{しゅう|き}が1\R{秒}{びょう}の\R{振}{ふ}り\R{子}{こ}の\R{糸}{いと}の\R{長}{なが}さは、\R{約}{やく}\qty{25}{cm}です。\R{周期}{しゅう|き}を2\R{秒}{びょう}にするには、\R{糸}{いと}の\R{長}{なが}さを\R{何倍}{なん|ばい}、\R{約}{やく}\R{何}{なん}cmにすればよいですか（2-8。\pageref{bk:2-15:bai}ページの\R{表}{ひょう}を\R{見}{み}よう）。
\end{minipage}\hfill
\begin{minipage}[t]{0.38\linewidth}\vspace{0mm}\centering
\begin{tikzpicture}[x=1mm, y=1mm]
  \foreach \cx/\L/\lab/\T in {8/7/{\R{約}{やく}\qty{25}{cm}}/{1\R{秒}{びょう}}, 38/28/{？}/{2\R{秒}{びょう}}} {
    \fill[bkGray!35] (\cx-8,0) rectangle (\cx+8,2.2);
    \draw[bkGray!90!black, line width=0.8pt] (\cx-8,0) -- (\cx+8,0);
    \draw[bkInk!60, line width=0.4pt, densely dashed] (\cx,0) -- (\cx,-\L-3);
    \draw[bkInk, line width=0.7pt] (\cx,0) -- ++(-75:\L);
    \shade[ball color=bkOrange] ($(\cx,0)+(-75:\L)+(-75:2.2)$) circle (2.2);
    \node[bk label, anchor=north] at (\cx,-\L-5) {\R{周期}{しゅう|き}\T};
    \node[font=\sffamily\footnotesize, text=bkInk, anchor=west] at ($(\cx,0)+(-75:\L/2)+(1.5,0)$) {\lab};}
\end{tikzpicture}
\end{minipage}
\bkpans{4\R{倍}{ばい}。\R{約}{やく}\qty{100}{cm}（\R{約}{やく}\qty{1}{m}）}

\bkproblem \qty{2}{N}で\R{引}{ひ}くと\qty{4}{cm}のびるばねがあります。このばねのばね\R{定数}{てい|すう}は\R{何}{なん}N/mですか。\qty{3}{N}で\R{引}{ひ}くと\R{何}{なん}cmのびますか。
\bkpans{$\qty{4}{cm}\div\qty{100}{cm/m}=\qty{0.04}{m}$、$k=\dfrac{2}{0.04}=50$ で\qty{50}{N/m}。$x=\dfrac{3}{50}=0.06$ で\qty{0.06}{m}、$\qty{0.06}{m}\times\qty{100}{cm/m}=\qty{6}{cm}$}

\bkproblem \R{乾電池}{かん|でん|ち}2\R{個}{こ}（\qty{3}{V}）につないだ\R{電熱線}{でん|ねつ|せん}に\qty{0.6}{A}\R{流}{なが}れました。この\R{電熱線}{でん|ねつ|せん}の\R{抵抗}{てい|こう}は\R{何}{なん}Ωですか。\R{乾電池}{かん|でん|ち}を4\R{個}{こ}（\qty{6}{V}）にすると、\R{電流}{でん|りゅう}は\R{何}{なん}Aになりますか。
\bkpans{$R=\dfrac{3}{0.6}=5$ で\qty{5}{\ohm}。\R{電圧}{でん|あつ}が2\R{倍}{ばい}なので\R{電流}{でん|りゅう}も2\R{倍}{ばい}で\qty{1.2}{A}}

\begin{furikaeri}
\begin{checklist}
\item \R{比例}{ひ|れい}を\R{見}{み}つけ、$y=ax$ の\R{式}{しき}とグラフにできる。\hfill{\footnotesize\sffamily\color{bkTheme}→\pageref{bk:2-15:hirei}ページ}
\item ばねの $F=kx$、オームの\R{法則}{ほう|そく} $V=RI$ を\R{使}{つか}える。\hfill{\footnotesize\sffamily\color{bkTheme}→\pageref{bk:2-15:bane}ページ}
\item \R{反比例}{はん|ぴ|れい}を\R{見分}{み|わ}け、2\R{倍}{ばい}にするとどうなるか\R{答}{こた}えられる。\hfill{\footnotesize\sffamily\color{bkTheme}→\pageref{bk:2-15:hanpirei}ページ}
\end{checklist}
\end{furikaeri}

\bkfukidashi{delta}{ばねも\R{電気}{でん|き}も、\R{表}{ひょう}を\R{作}{つく}ってわり\R{算}{ざん}をしたら、\R{同}{おな}じ「\R{比例}{ひ|れい}」の\R{形}{かたち}が\R{見}{み}えたね。\R{次}{つぎ}の\R{章}{しょう}では、\R{追}{お}いかけっこを\R{式}{しき}で\R{解}{と}くよ。}
